% ============================================================================
\section{Cache}
\label{sec:cache}

Cache organization, block replacement (LRU, FIFO, Tree-LRU, Clock), write
policies, the three categories of misses, and cache performance (AMAT, hit and miss
rates, multilevel caches).

% ----------------------------------------------------------------------------
\subsection{What to know}

\subsubsection{Basics}
\begin{itemize}[leftmargin=*]
  \item Memory hierarchy: each level is larger, slower and cheaper per byte than
  the one above. It works because of \textbf{locality}: \emph{temporal} (an item
  used now is used again soon) and \emph{spatial} (nearby items are used soon, so
  whole blocks are moved).
  \item \textbf{Hit}: the data is in the cache. \textbf{Hit time}: time to access
  the cache, including deciding hit or miss. \textbf{Miss penalty}: time to bring
  the block from the lower level. Hit rate $+$ miss rate $=1$.
  \item The four questions: \emph{placement} (where can a block go),
  \emph{identification} (how is it found: tag $+$ valid bit), \emph{replacement}
  (which block is evicted), \emph{write strategy}.
\end{itemize}

\subsubsection{Organization and address breakdown}
$C$ = cache size, $B$ = block size, $a$ = associativity (ways), $A$ = address bits.
\[
  \#\text{blocks}=\frac CB,\qquad \#\text{sets } S=\frac{C}{B\times a},
\]
\[
  \text{offset}=\log_2B,\quad \text{index}=\log_2S,\quad \text{tag}=A-\text{index}-\text{offset}.
\]
\begin{center}
\small
\begin{tabular}{@{}lll@{}}
\toprule
Organization & Where a block can go & Index bits \\
\midrule
Direct mapped ($a=1$) & exactly one block: (block address) mod (\#blocks) & $\log_2(C/B)$ \\
$n$-way set associative & any of the $n$ ways of set (block address) mod $S$ & $\log_2(C/(Bn))$ \\
Fully associative ($S=1$) & anywhere; compare all tags & 0 \\
\bottomrule
\end{tabular}
\end{center}
\begin{itemize}[leftmargin=*]
  \item Block address $=\lfloor\text{addr}/B\rfloor$; set $=$ block address mod $S$;
  tag $=\lfloor\text{block address}/S\rfloor$.
  \item Doubling associativity (same $C$, $B$): index $-1$ bit, tag $+1$ bit.
  \item Storage per block $=8B$ data bits $+$ tag $+$ valid ($+$ dirty for write-back)
  ($+$ replacement bits).
\end{itemize}

\subsubsection{Block replacement}
Only needed when $a>1$ (direct mapped has no choice).
\begin{itemize}[leftmargin=*]
  \item \textbf{Random}: pseudo-random way; simple, no state.
  \item \textbf{LRU}: evict the block unused for the longest time. Exact LRU
  keeps the full use order of the set ($\lceil\log_2 n!\rceil$ bits, e.g.\ 5 bits
  for 4 ways), which is costly for large $n$.
  \item \textbf{FIFO}: evict the block loaded earliest; hits do not change the
  order.
  \item \textbf{Tree-LRU} (tree pseudo-LRU): $n-1$ bits in a binary tree over the
  ways (3 bits for 4 ways: root $b_0$ chooses a pair, $b_1$ the way in pair
  $\{0,1\}$, $b_2$ the way in pair $\{2,3\}$). On every access, set the bits on the
  path to point \emph{away} from the accessed way. Victim: follow the bits from the
  root. It approximates LRU and can evict a block that true LRU would keep.
  \item \textbf{Clock} (second chance): ways in a circle, one reference bit each and a
  pointer. On an access set the bit to 1. To replace: if the pointed block has bit 1,
  clear it and advance; evict the first block with bit 0. The new block gets bit 1
  and the pointer moves past it.
\end{itemize}

\subsubsection{Write policies}
\begin{center}
\small
\begin{tabular}{@{}p{2.6cm}p{11.6cm}@{}}
\toprule
Write-through & on a write hit, update the cache \emph{and} the next level; no dirty
bit; a write buffer hides the stalls; memory is always up to date. \\
Write-back & update only the cache; a \textbf{dirty bit} marks modified blocks, which
are written to memory only when evicted; less memory traffic. \\
\midrule
Write allocate & on a write miss, fetch the block into the cache, then write it. \\
No-write allocate & on a write miss, write only to the next level; the cache is not
changed. \\
\bottomrule
\end{tabular}
\end{center}
Usual pairs: write-through $+$ no-write allocate; write-back $+$ write allocate.

\subsubsection{The three categories of misses (3 Cs)}
\begin{itemize}[leftmargin=*]
  \item \textbf{Compulsory} (cold start): the first access to a block.
  \item \textbf{Capacity}: the cache cannot hold all the blocks the program needs;
  the miss also happens in a fully associative LRU cache of the same size.
  \item \textbf{Conflict} (collision): too many blocks map to the same set; the
  access would hit in a fully associative cache of the same size.
  \item (A fourth C, \textbf{coherence}, appears in multiprocessors.)
\end{itemize}
\begin{center}
\small
\begin{tabular}{@{}p{4.2cm}p{4.6cm}p{5.4cm}@{}}
\toprule
Basic optimization & Reduces & Cost \\
\midrule
Larger block size & compulsory misses (spatial locality) & larger miss penalty; more conflict/capacity misses if too large \\
Larger cache & capacity misses & longer hit time, cost, power \\
Higher associativity & conflict misses & longer hit time \\
Multilevel caches & miss penalty & more hardware \\
Read misses before writes & miss penalty & check the write buffer \\
Avoid translation in indexing & hit time & needs VIPT, $C/a\le$ page size \\
\bottomrule
\end{tabular}
\end{center}
Rule of thumb: a direct-mapped cache of size $N$ misses about as often as a 2-way
cache of size $N/2$.

\subsubsection{Cache performance}
\[
  \text{AMAT}=\text{Hit time}+\text{Miss rate}\times\text{Miss penalty}
\]
\[
  \CPI=\CPI_{exec}+\frac{\text{memory accesses}}{\text{instruction}}\times\text{Miss rate}\times\text{Miss penalty},
  \qquad T=I\times\CPI\times C.
\]
Accesses per instruction $=1$ (fetch) $+$ fraction of loads/stores. Miss penalty
in \emph{cycles} $=$ penalty in ns / clock period.

\textbf{Multilevel caches} (L2 reduces the L1 \emph{miss penalty}):
\[
  \text{AMAT}=\text{HT}_{L1}+\text{MR}_{L1}\times\big(\text{HT}_{L2}+\text{MR}_{L2}\times\text{MP}_{L2}\big)
\]
\begin{itemize}[leftmargin=*]
  \item \textbf{Local} miss rate $=$ misses in this cache / accesses \emph{to this
  cache} (for L2: accesses $=$ L1 misses). \textbf{Global} miss rate $=$ misses in
  this cache / all CPU accesses $=\text{MR}_{L1}\times\text{MR}_{L2,local}$.
  \item Stalls per instruction $=\dfrac{\text{L1 misses}}{\text{instr}}\times
  \text{HT}_{L2}+\dfrac{\text{L2 misses}}{\text{instr}}\times\text{MP}_{L2}$.
\end{itemize}

% ----------------------------------------------------------------------------
\subsection{Common traps}
\begin{enumerate}[leftmargin=*]
  \item Sets $=C/(B\times a)$: divide by the associativity.
  \item Work with block addresses (addr$/B$), not byte addresses, when tracing.
  \item L2 \emph{local} miss rate is relative to L1 misses; multiply by
  $\text{MR}_{L1}$ for the global rate.
  \item Count the instruction fetch in memory accesses per instruction.
  \item Tree-LRU is not LRU: trace the bits, do not guess.
  \item In write-back, a miss that evicts a dirty block costs a write-back as well.
  \item A conflict miss is one that a fully associative cache of the same size
  would \emph{not} have had.
\end{enumerate}

% ----------------------------------------------------------------------------
\subsection{Lecturer's review questions}

\begin{exercise}[Lecturer 1]
The algorithm that removes contents and places new contents into the cache is
called: (a) replacement algorithm; (b) renewal algorithm; (c) updation; (d) none of
the above.
\end{exercise}
\begin{solution}
\textbf{(a)}.
\end{solution}

\begin{exercise}[Lecturer 2]
An L2 cache helps to: (a) reduce the miss rate of L1; (b) reduce the hit time of
L1; (c) double the L1 size; (d) reduce the miss penalty of L1.
\end{exercise}
\begin{solution}
\textbf{(d)}: an L1 miss is served by L2 (a few cycles) instead of memory (about
100 cycles). L1's own miss rate and hit time do not change.
\end{solution}

\begin{exercise}[Lecturer 8--11]\label{l:cache-trace}
A CPU has an 8~KB, 2-way set-associative data cache with 16-byte lines and 32-bit
addresses.
(8) Offset bits: (a) 2; (b) 4; (c) 8; (d) 16.
(9) Index bits: (a) 2; (b) 4; (c) 8; (d) 16.
(10) Tag bits: (a) 16; (b) 18; (c) 20; (d) 22.
(11) After reset, how many misses for \texttt{0x00000000}, \texttt{0x00000014},
\texttt{0x00000008}, \texttt{0x00000010}, \texttt{0x00000020}? (a) 2; (b) 3; (c) 4;
(d) 5.
\end{exercise}
\begin{solution}
(8) \textbf{(b)}: $\log_216=4$. (9) \textbf{(c)}: $S=\dfrac{8192}{16\times2}=256$,
$\log_2256=8$. (10) \textbf{(c)}: $32-8-4=20$.

(11) \textbf{(b) 3}. Block address $=$ addr$/16$:
\begin{center}
\small
\begin{tabular}{@{}lccccc@{}}
\toprule
Address & \texttt{0x00} & \texttt{0x14} & \texttt{0x08} & \texttt{0x10} & \texttt{0x20} \\
Block (set) & 0 & 1 & 0 & 1 & 2 \\
Result & miss & miss & hit & hit & miss \\
\bottomrule
\end{tabular}
\end{center}
\texttt{0x08} is in the same 16-byte block as \texttt{0x00}, and \texttt{0x10} in the
same block as \texttt{0x14} (spatial locality). All three misses are compulsory.
\end{solution}

% ----------------------------------------------------------------------------
\subsection{Exam-style questions}
Solutions are in Section~\ref{sec:cache-sol}.

\begin{exercise}[Address breakdown]\label{q:c-addr}
A 32~KiB cache has 64-byte blocks and 32-bit byte addresses.
\begin{enumerate}[label=(\alph*),leftmargin=*]
  \item Give the tag, index and offset widths if it is direct mapped, 4-way set
  associative, and fully associative.
  \item For the 4-way write-back version, find the total storage (data, tag, valid
  and dirty bits).
  \item In the 4-way cache, which set and tag does address \texttt{0x12345678} map
  to?
\end{enumerate}
\end{exercise}

\begin{exercise}[Replacement policies]\label{q:c-repl}
One set of a 4-way set-associative cache, initially empty, receives these block
references:
\begin{center}
\texttt{0 1 2 3 0 4 1 0 5 2 3 0}
\end{center}
\begin{enumerate}[label=(\alph*),leftmargin=*]
  \item Trace LRU, FIFO, Tree-LRU and Clock. Mark hits and misses and give the
  victim of each miss. (Fill empty ways in order 0--3. Tree-LRU bits are
  $b_0b_1b_2$, initially 000; Clock gives new blocks bit 1.)
  \item Count the misses of each policy. Where does Tree-LRU differ from LRU?
  \item How many state bits per set does each policy need?
\end{enumerate}
\end{exercise}

\begin{exercise}[The 3 Cs]\label{q:c-3c}
A direct-mapped cache has 4 one-word blocks (index $=$ word address mod 4) and
starts empty. Classify every access to the word addresses
\texttt{0 4 0 1 2 3 5 0 4 1} as hit, compulsory, capacity or conflict miss.
\end{exercise}

\begin{exercise}[Write policies]\label{q:c-write}
A direct-mapped cache of 4 one-word blocks starts empty. The CPU performs
\texttt{W0 R0 W4 W4 R0 W8 R8 W0} (W $=$ write, R $=$ read, word addresses). For
(a) write-through with no-write allocate and (b) write-back with write allocate,
count the misses, the blocks read from memory and the words written to memory.
Which blocks are dirty at the end of (b)?
\end{exercise}

\begin{exercise}[Multilevel cache]\label{q:c-amat}
L1 hit time 1 cycle, L1 miss rate 2\%; L2 hit time 10 cycles, L2 local miss rate
40\%; memory 100 cycles. There are 1.3 memory accesses per instruction and
$\CPI_{exec}=1$.
\begin{enumerate}[label=(\alph*),leftmargin=*]
  \item Find the AMAT with and without L2, and the L2 global miss rate.
  \item Find the CPI with and without L2 and the speedup due to L2.
  \item Which of the 3 Cs does L2 affect? Which cache parameter does it improve?
\end{enumerate}
\end{exercise}

% ----------------------------------------------------------------------------
\subsection{Solutions}
\label{sec:cache-sol}

\begin{noot}{Exercise~\ref{q:c-addr}}
\begin{enumerate}[label=(\alph*),leftmargin=*]
  \item Offset $=\log_264=6$; $C/B=512$ blocks.
  Direct mapped: index 9, tag $32-9-6=17$.
  4-way: $S=128$, index 7, tag 19.
  Fully associative: index 0, tag 26.
  \item Per block: $512$ data $+19$ tag $+1$ valid $+1$ dirty $=533$ bits;
  $512\times533=272{,}896$ bits $=34{,}112$ bytes (about 4\% overhead).
  \item \texttt{0x12345678}$/64=$ \texttt{0x48D159}; set $=$ \texttt{0x48D159} mod
  128 $=$ \texttt{0x59} $=89$; tag $=$ \texttt{0x48D159}$/128=$ \texttt{0x91A2}.
\end{enumerate}
\end{noot}

\begin{noot}{Exercise~\ref{q:c-repl}}
(a) M $=$ miss (victim in brackets), H $=$ hit.
\begin{center}
\small
\setlength{\tabcolsep}{3pt}
\begin{tabular}{@{}lcccccccccccc@{}}
\toprule
Ref & 0 & 1 & 2 & 3 & 0 & 4 & 1 & 0 & 5 & 2 & 3 & 0 \\
\midrule
LRU & M & M & M & M & H & M(1) & M(2) & H & M(3) & M(4) & M(1) & H \\
FIFO & M & M & M & M & H & M(0) & H & M(1) & M(2) & M(3) & M(4) & H \\
Tree-LRU & M & M & M & M & H & M(2) & H & H & M(3) & M(1) & M(4) & H \\
\quad bits $b_0b_1b_2$ & 110 & 100 & 001 & 000 & 110 & 011 & 101 & 111 & 010 & 100 & 001 & 111 \\
Clock & M & M & M & M & H & M(0) & H & M(2) & M(3) & M(1) & M(4) & H \\
\bottomrule
\end{tabular}
\end{center}
Tree-LRU detail at reference 4: the bits are 110 ($b_0=1$: victim in pair
$\{2,3\}$; $b_2=0$: way 2), so block 2 is evicted, while true LRU evicts 1. Clock at
reference 4: all bits are 1, so the pointer clears all four and comes back to way 0,
evicting 0.

(b) LRU 9, FIFO 9, Tree-LRU 8, Clock 9. Tree-LRU differs at reference 4 (it keeps 1,
which is reused next) and from then on holds different blocks. Here the
approximation happens to do better; in general it is slightly worse than LRU.

(c) Exact LRU: $\lceil\log_2 4!\rceil=5$ bits (often 8: a 2-bit age per way).
FIFO: a 2-bit pointer. Tree-LRU: 3 bits. Clock: 4 reference bits $+$ a 2-bit
pointer. Random: none.
\end{noot}

\begin{noot}{Exercise~\ref{q:c-3c}}
\begin{center}
\small
\begin{tabular}{@{}lcccccccccc@{}}
\toprule
Address & 0 & 4 & 0 & 1 & 2 & 3 & 5 & 0 & 4 & 1 \\
Index & 0 & 0 & 0 & 1 & 2 & 3 & 1 & 0 & 0 & 1 \\
Result & comp. & comp. & confl. & comp. & comp. & comp. & comp. & hit & cap. & cap. \\
\bottomrule
\end{tabular}
\end{center}
The third access (0) misses only because 4 took index 0; a 4-block fully
associative LRU cache would still hold 0, so it is a conflict miss. The last two
(4 and 1) also miss in the fully associative cache, which has meanwhile evicted them
(too many distinct blocks since their last use), so they are capacity misses.
\end{noot}

\begin{noot}{Exercise~\ref{q:c-write}}
\begin{enumerate}[label=(\alph*),leftmargin=*]
  \item Write-through, no-write allocate: writes never fill the cache. W0 miss, R0
  miss (fetch 0), W4 miss, W4 miss, R0 hit, W8 miss, R8 miss (fetch 8, replaces 0),
  W0 miss. \textbf{7 misses, 2 blocks read, 5 words written} (every write goes to
  memory).
  \item Write-back, write allocate: W0 miss (fetch 0, dirty), R0 hit, W4 miss
  (write back 0, fetch 4, dirty), W4 hit, R0 miss (write back 4, fetch 0, clean), W8
  miss (0 is clean: no write-back; fetch 8, dirty), R8 hit, W0 miss (write back 8,
  fetch 0, dirty). \textbf{5 misses, 5 blocks read, 3 words written}; block 0 is
  dirty at the end (one more write-back when it is evicted).
\end{enumerate}
\end{noot}

\begin{noot}{Exercise~\ref{q:c-amat}}
\begin{enumerate}[label=(\alph*),leftmargin=*]
  \item With L2: $1+0.02\times(10+0.4\times100)=1+0.02\times50=2$ cycles. Without:
  $1+0.02\times100=3$ cycles. Global L2 miss rate $=0.02\times0.4=0.8\%$.
  \item With L2: stalls $=1.3\times0.02\times50=1.3$, $\CPI=2.3$. Without:
  $1.3\times0.02\times100=2.6$, $\CPI=3.6$. Speedup $=3.6/2.3=1.57$.
  \item L2 does not change L1's misses (any C); it reduces the \emph{miss penalty}
  of L1, because most L1 misses now hit in L2.
\end{enumerate}
\end{noot}

